\section{Methodology}

\subsection{Study Design and Evidence Window}

This study performs a descriptive, build-aligned analysis of 15 Android consumer apps. It examines three analytical dimensions: the static exposure present in selected app builds, the runtime behavior observed for those builds, and the relationships or divergences between the two evidence layers. The unit of inference is the selected app-build group, defined by package name, version code and name, base APK hash, split-APK identity, contributing static run, and selected dynamic runs.

Collection was conducted on a non-root Motorola Moto G 5G (2024) running Android 15. Evidence was selected from a fixed window spanning June 26 through July 9, 2026 UTC, used as a recency, build-alignment, and provenance control rather than as a longitudinal repeatability experiment. Each app was grouped by primary user-facing function into News, Social Media, Messaging, and Professional Networking categories. The categories are used for descriptive visualization only and are not treated as statistical strata. The non-root setting avoided app modification, privileged instrumentation, TLS interception, and payload inspection.

\subsection{Static and Runtime Evidence}

Static analysis used the harvested base and split APK artifacts associated with each selected build. The reported variables include dangerous permissions, high- and medium-severity findings, unguarded-export findings, network- and storage-related findings, privacy and platform findings, MASVS-aligned categories, and supported SDK or service indicators \cite{owaspMasvs2024}. Unguarded-export findings are High- or Medium-severity detector rows whose titles indicate a missing permission or a weak guard; they count finding rows, not unique Android components. These measurements represent descriptive exposure indicators rather than calibrated exploitability, compliance, or risk scores.

Runtime analysis used package-associated evidence packs and PCAP-derived metadata. Captures were classified as strict-idle, quiescent foreground (QFG), or interactive. Strict-idle captures represent quiet foreground baselines; QFG captures represent no-touch foreground periods in which app-driven activity continued; and interactive captures exercise user-facing behavior. Runtime variables include capture count and duration, PCAP volume, packet rate, retained DNS/SNI hostname breadth, service families, protocol prevalence, and baseline-relative changes where comparable baseline and interactive evidence were available. Hostname breadth is the per-run union of retained top DNS and TLS-SNI names, not an exhaustive inventory of endpoints. Retained interactive DNS/SNI hostnames were matched by exact or subdomain relation against tracking-oriented categories in a pinned Disconnect services list; each app's documented own registrable domain (eTLD+1) was excluded from the third-party count \cite{disconnectTrackingProtection}.

\subsection{Build Selection and Evidence Alignment}

Qualification policy $P$ is the historical technical-validity rule used for the submitted cohort. A dynamic run qualifies under $P$ when it has a nonempty run identifier, \texttt{valid\_dataset\_run} status, required protocol metadata (profile, sequence, and interaction level), duration of at least 180\,s, at least 20 analysis windows of width 10\,s and stride 5\,s, a PCAP of at least 50\,000\,B (10\,000\,B for qualifying connected or text-messaging captures), and a usable PCAP analysis or an explicit no-traffic result.

Qualifying runs are grouped by package, version code and name, and base APK hash. One build group per app is then chosen by an ordered ranking on readiness, coverage of evidence classes, capped run counts, PCAP and QA state, recency, match to the then-current install, and version code. Static and runtime observations are integrated only when they refer to the same selected build. Evidence with mismatched identity, missing support, invalid QA state, or insufficient signal is excluded or deferred. Algorithm~\ref{alg:build-aligned-selection} summarizes this procedure.

\begin{algorithm}[!t]
\caption{Build-Aligned Static-Runtime Evidence Selection}
\label{alg:build-aligned-selection}
\begin{algorithmic}[1]
\REQUIRE cohort $A$, 14-day window $W$, static set $S$, dynamic set $D$, qualification policy $P$
\ENSURE aligned app-build rows, excluded/deferred rows, source data, checksums
\FOR{each app $a \in A$}
  \STATE retrieve candidate dynamic runs for $a$ within $W$
  \STATE retain runs satisfying technical qualification policy $P$
  \STATE group runs by normalized package, version code/name, and base APK SHA-256
  \STATE summarize evidence-class coverage and readiness for each group
  \STATE select one group using the ordered build-ranking policy
  \STATE resolve static evidence in $S$ for the selected build
  \IF{static and dynamic identities do not align}
    \STATE exclude or defer $a$; persist the reason; continue
  \ENDIF
  \STATE preserve strict-idle, QFG, and interactive classes separately
  \STATE compute static exposure and runtime behavior summaries
  \STATE add the aligned app-build row and provenance
\ENDFOR
\STATE compute descriptive cross-layer profiles and app-level associations
\STATE export aligned rows, excluded/deferred rows, source data, and checksums
\end{algorithmic}
\end{algorithm}

\subsection{Integrated Analysis and Reproducibility}

\emph{Risk posture} denotes evidence-supported conditions that may create or enlarge opportunities for privacy or security harm; it is not a calibrated probability of exploitation or realized harm. Static indicators characterize potential packaged exposure; runtime indicators characterize observed behavior. The two layers are interpreted jointly and are not collapsed into a single score.

The study reports descriptive counts, evidence-class medians, baseline-relative packet-rate changes, and exploratory Spearman associations. The reported median packet-rate shift is the app-level median of $\log_2[(I+1)/(B+1)]$, where $I$ is the interactive median packets/s and $B$ is the strict-idle median if that class is present and the QFG median otherwise. In this cohort, only TikTok lacks a selected strict-idle capture and therefore uses QFG as the baseline $B$. Cross-layer profiles use a transparent median split: higher static exposure is high- and medium-severity finding volume at or above the cohort median (52), and higher runtime activity is interactive retained-hostname count at or above the cohort median (12). Equality is assigned to the higher group. These profiles are cohort-relative and are not a composite risk score or predictive model.

PCAP processing is limited to metadata and derived transport indicators. Message bodies, cookies, authorization headers, private account content, and decrypted TLS payloads are not inspected. Reproducibility artifacts include the selected-build manifest, source CSVs, generated tables and figures, inclusion and exclusion records, SHA-256 checksums, and the pinned Disconnect list used for hostname association.

\subsection{Limitations}

The study is limited to 15 apps, one Android 15 physical device, one collection environment, selected interaction scenarios, and the available static detectors. Device and environment consistency improve within-study comparability but limit generalization across OEM implementations, Android versions, hardware, and networks. Runtime observations are scenario-bounded and do not establish exhaustive app behavior, while unequal capture depth may affect the precision of app-level comparisons. Multiple captures characterize within-window behavior for the selected builds; the single evidence window does not establish temporal repeatability across independent collection periods. The non-root, no-MITM design preserves a clear privacy boundary but prevents inspection of encrypted payload semantics. Accordingly, elevated packet rate or transfer volume may reflect media delivery, synchronization, telemetry, advertising, keepalive traffic, or privacy-sensitive exchange; transport metadata alone cannot distinguish among these explanations. A Disconnect blocklist match indicates association with categorized infrastructure; it does not reveal encrypted payload contents or establish private-data leakage.

Static finding counts are tool-dependent exposure indicators rather than formal measures of exploitability, MASVS compliance, or privacy harm. Finding volume is exposure volume, not a severity-weighted realized-risk score. Observed runtime activity does not demonstrate malicious behavior, and limited activity does not establish low risk. The analysis therefore characterizes the static and runtime risk posture of the selected builds without inferring causation, predicting future behavior, or generalizing the results to all Android devices and app versions.
